
We treat AlphaFold3 not as a second benchmark study but as a portability
question: does what Section~\ref{sec:parallelism} and
Section~\ref{sec:results-hardware} found about AlphaFold2 on TPU
generalize to a newer model in the same family. AlphaFold3 is run with
its real, trained weights, not the untrained weights used for AlphaFold2
elsewhere in this paper (Section~\ref{sec:methodology}). The setup,
timings and per-sample score comparisons are in
Appendix~\ref{app:af3}; two findings bear on the argument of this
paper.

\subsection{No TPU backend in the release we tested}
\label{sec:af3-tpu}

Running the same TPU Job configuration used throughout this paper
against AlphaFold3 fails immediately: \texttt{run\_alphafold.py} rejects
\texttt{-{}-jax\_backend=tpu} at flag parsing, after 5.588\,s, before any
model code executes. The accepted values are \texttt{cpu}, \texttt{gpu}
and \texttt{mps}.
This matches AlphaFold3's own documentation, which lists CPU or an
NVIDIA GPU of compute capability 7.0 or greater as the supported
targets \citep{af3docs2026}. TPU is not a documented target. The
rejection above is a Job failure on our own cluster, not an inference
from the documentation alone.


\subsection{Output differs across backends}
\label{sec:af3-body-repro}

Running the identical seed and input on the same backend but different
machines (Stanford's CPU cluster versus the Colab CPU) changes the
per-sample ranking score by at most 1.67\%, and by under 0.1\% on four
of five samples; running it on different backends (Colab CPU versus
Colab GPU) changes it by up to 32.33\% on one sample.
The across-backend divergence is roughly an order of magnitude larger
than the across-machine, same-backend noise. A plausible, documented
cause exists in AlphaFold3's own performance notes, which report known
numerical issues on CUDA Capability 7.x devices such as the Colab T4
\citep{af3docs2026}. We did not isolate the operation responsible, so
this cause remains plausible and unconfirmed
(Appendix~\ref{app:af3-repro}).

\subsection{Summary}
\label{sec:af3-summary}

Section~\ref{sec:parallelism}'s central finding, that realized
performance on TPU depends on how parallelism is expressed rather than
on the accelerator alone, cannot be tested on AlphaFold3 at all: the
release we tested has no TPU path to express parallelism on in the
first place. The broader caution of this paper does carry over: a
newer model in the same family changes not just speed but which
backends are available at all, and, on a documented edge case of GPU
hardware, how far the reported output can be trusted to match across
sessions.
